% ============================================================
% DRAFT — corrected "Plan correctness under vague commands" (RQ1)
% section, to replace the current \mysection{...}{...}
% \label{sec:rq1_results} ... block in evaluation.tex (currently
% the block running from "Table~\ref{tab:rq1_levels} summarises..."
% down to the paragraph ending "...dialogue state carry-over is
% fixed." right before \mysection{Executability on the Physical
% Chassis}).
%
% This file is a draft for review, not a finished chapter section.
% Read it next to the current evaluation.tex, copy over whatever
% you're happy with, and edit the rest — nothing here has been
% written into the actual chapter file.
%
% What changed from the current chapter text, and why:
%   1. The table now covers all 41 tagged physical runs, not 23.
%      The other 18 were already sitting in the logs with a
%      test_level set; they just weren't in the table.
%   2. Three L0 "describe scene" rows are actually named-target
%      search commands ("go to the orange box", "go to the orange
%      box in the room") or a guidance-loop continuation utterance
%      ("can you please continue the plan") and move to L3.
%      Three more L0 rows combine a named target with a relative
%      direction ("on your left", "to your right") and move to L5.
%      This is the only recoding done; nothing else changed level.
%   3. The table caption said Gemma 4 E4B was used throughout. Of
%      the 41 runs, 34 used Gemma for both roles and 7 used
%      Qwen2.5-VL 3B as the vision model, so the caption is
%      corrected.
%   4. A second pair of columns reports the operator's own
%      good/bad call from run_flags.json next to the pipeline's
%      outcome label, because the two disagree often enough to be
%      a finding on their own.
%   5. The per-level paragraphs are rewritten against the
%      corrected counts and quote the operator notes that actually
%      explain the halts, rather than leaving candidatetodo
%      placeholders.
% ============================================================

\mysection{Plan correctness under vague commands}{Plan Correctness under Vague Commands}
\label{sec:rq1_results}

Table~\ref{tab:rq1_levels} covers all 41 physical runs that carry a command-level tag in the logs, recorded between
8 and 11 October: 34 with Gemma~4 E4B as both planner and vision model, 7 with Qwen2.5-VL~3B as the vision model.
The previous version of this table used only 23 of these runs and misfiled six of them. Three commands that name a
target directly (``go to the orange box'', ``go to the orange box in the room'', and ``can you please continue the
plan'', which resumes an approach already in progress) were logged as level~0 but are perception-dependent
searches, so I moved them to level~3. Three more combine a named target with a relative direction (``go to the
black bag on your left'', ``\ldots{} the orange box on the right of the black bag'', ``\ldots{} to the orange box
to your right'') and belong at level~5, not level~0. No other run changed level. I also added a column for the
operator's own good/bad call from \texttt{run\_flags.json}, next to the pipeline's own outcome label, because the
two do not always agree, and where they disagree is informative.

\begin{table}[htbp]
    \centering
    \caption{Physical rover outcomes per command level, all 41 tagged runs (34 Gemma~4 E4B, 7 Qwen2.5-VL~3B for vision), against the pipeline's own outcome label and the operator's recorded verdict.}
    \label{tab:rq1_levels}
    \small
    \setlength{\tabcolsep}{4.5pt}
    \begin{tabular}{lp{2.9cm}rrrrrr}
        \toprule
        Level & Type & Runs & Done & After guidance & Halted/ aborted & Operator decision good & Operator decision bad \\
        \midrule
        L0 & Describe scene (control) & 6 & 5 & 0 & 1 & 4 & 0 \\
        L1 & Single action & 4 & 4 & 0 & 0 & 4 & 0 \\
        L2 & Sequence of actions & 4 & 4 & 0 & 0 & 4 & 0 \\
        L3 & Perception-dependent & 8 & 3 & 2 & 3 & 5 & 1 \\
        L4 & Conditional & 9 & 7 & 1 & 1 & 4 & 2 \\
        L5 & Vague & 10 & 3 & 3 & 4 & 5 & 2 \\
        \midrule
        & Total & 41 & 26 & 6 & 9 & 26 & 5 \\
        \bottomrule
    \end{tabular}
\end{table}
\candidatetodo{10 of the 41 runs have no operator flag at all (good+bad only add to 31). }

\paragraph{Levels 1 and 2.} All eight runs produced the expected plan on the first attempt and finished without
intervention. These levels need no perception, so they show the planner maps simple action chains onto the
capability vocabulary without help, though eight runs is still too few to call it a rate. The runs did surface a
hardware issue worth flagging here even though it sits outside plan correctness: in the earliest L1/L2 runs the
rover turned about 118\textdegree{} where 90\textdegree{} was commanded, which the operator noted happened ``even
though I calibrated it''. The turn timing was retuned partway through testing (one note records changing it to
0.76\,s), and later runs in the same batch turn at close to 90\textdegree{}. \candidatetodo{This is an open-loop
timing calibration issue, not a planning or vision one, so it probably belongs in Section~\ref{sec:executability}
next to the sonar and distance results rather than here. Move it there, and say whether the 0.76\,s value is what
the final build uses.}

\paragraph{Level 0.} Of the six genuine describe-scene runs, five completed and one (``so tell me what you can
see'') halted after guidance with no operator note explaining why. \candidatetodo{Check this run's log directly
(session 2d9a17236f51) before writing a cause for it; I don't want to guess.}

\paragraph{Level 3.} Of the eight runs, three completed directly, two completed after a guidance turn, and three
halted or aborted. Two of the three failures trace to the room's lighting rather than to planning: in session
2824ea8048bb (``please don't do the black bag'', which the pipeline resolved to driving at the bag — the negation
loss already noted in earlier drafts of this chapter) the operator wrote that the model ``locked onto another black
object since it is hard to distinguish a black bag in the dark room''. In session caf25eb230ef the rover found the
box, drove into an obstacle on the way, and reported the mission complete anyway because the box was still visible
in the background frame — the arrival check accepted a frame that happened to contain the target without checking
that the rover had actually stopped in front of it. That one is a genuine false positive: the operator flagged it
``review'', not ``good'', for exactly this reason. The third, session f13aa9ac3c8d, is more a success than a
failure: the rover lost the box, did a 360\textdegree{} search, got three guidance turns, relocated it, and drove
there, and only ended on \texttt{guidance\_exhausted} because the arrival check could not confirm the final frame
(the operator's note: the box ``blended into the cupboard door behind it''). The operator flagged this run ``good''
despite the halted outcome, which is the clearest case in the batch of the pipeline's own outcome label and the
operator's judgement disagreeing: the guidance loop did what it was supposed to do, and the pipeline's label still
reads as a failure.

\paragraph{Level 4.} Seven of nine runs completed outright, one completed after guidance, and one aborted at an
obstacle. The two operator-flagged failures are a dropped condition and a misidentified target rather than a
hallucinated plan: in session abe08c4094f1 the operator's note is simply ``needed to check for condition'', meaning
the \texttt{check} step the command asked for was skipped; in session 1dddfd8d1f6a the rover ``went to the wrong
object and thought a jersey was the black bag''. A third run (f8811861d4e0, flagged ``review'') reported completion
with almost no information back, which the operator recorded as ``it did nothing, could not view correctly''. None
of these three is the condition itself failing to resolve correctly, which is closer to what I'd expect to be the
hard part of level~4 — they are a skipped step, a vision misclassification, and a near-empty report.

\paragraph{Level 5.} Ten runs, the most of any level and the most mixed: three completed, three completed after
guidance, four halted or aborted. Lighting shows up here again and more directly than at level~3. In session
2acdffb1c7c1 the operator wrote that the model ``kept skipping the frame as it cannot actually see the box due to
the light in the room not being enough'', and in session 794bd92f55de (``go to the black bag and then go to the
orange box after'') the rover reached the bag, then stalled at the box in the dark and only succeeded once the
operator ``moved it into the light and it could see it better'' — a direct, causal demonstration that the failure
was illumination, not planning, because moving the same object into better light fixed it mid-run. Session
26285fd30325, one of the three reclassified into this level, halted for the same reason without a note explaining
it. Session 0df36e214a46 completed but under-reported, telling the operator only that it had arrived rather than
what it found there. Session e39230867caa shows the clarification loop working as intended on a badly garbled
transcript (``go towards the rn fox and then come to see a rally'' resolved to ``go to the orange box and tell me
what you see around it''), which is the one unambiguous success story at this level. Taken together, level~5's
problems are less about the planner inventing structure for a vague command — the plans themselves were reasonable
in every run I read — and more about the vision model running out of light to work with once the rover is close
enough to need a fine judgement.

\paragraph{Lighting as the dominant failure cause.} Across levels~3 and~5, every halt with an operator note
attached to it names poor lighting or low contrast as the proximate cause: a black bag against a dark floor, an
orange box blending into a cupboard door, a frame rejected as ``too blurry'' because the room was too dim, and one
case where moving the object into better light fixed the run on the spot. The one clear exception is the
obstacle-collision-reported-as-complete case at level~3, which is an arrival-check problem rather than a perception
one. \candidatetodo{This is worth a sentence in the Limitations section too: the test room's lighting wasn't
controlled or measured, so the 41 runs are really a test of the pipeline under whatever light happened to be in the
room that day, not a controlled low-light condition.}

\paragraph{Where the pipeline's own label and the operator's verdict disagree.} Five runs are flagged ``bad'' and
five more ``review'', for ten flagged problems on top of (not always matching) the nine pipeline-labelled halts and
aborts. The two don't line up one-to-one: session f13aa9ac3c8d is labelled \texttt{halted\_after\_guidance} by the
pipeline but ``good'' by the operator, because the recovery behaviour was correct even though the mission didn't
finish; session caf25eb230ef is labelled \texttt{completed} by the pipeline but ``review'' by the operator, because
the mission finished on a technicality the arrival check should have caught. A pipeline-outcome count alone would
have reported this batch as 32 of 41 successful (26 completed plus 6 after guidance); the operator's count says 26
good against 5 explicit bad, with 10 runs never reviewed at all. \candidatetodo{Decide which number is the headline
result and say why. I'd lean toward reporting both, the way this paragraph does, rather than picking one — the gap
between them is itself a finding about how much the raw outcome label can be trusted.}

\paragraph{Exploratory runs.} Earlier runs between 7 and 9 October used different model pairs and are not part of
Table~\ref{tab:rq1_levels} (these predate the command-level tagging, which is why they have no \texttt{test\_level}
in the logs). In six search-and-approach runs with Gemma~4 E4B as the vision model, five halted after guidance with
the target not found and one completed. In seven search-and-approach runs with Qwen2.5-VL~3B as the vision model,
three completed directly, two completed after guidance, and two aborted. Six further runs on 9 October (levels 0 to
4, Qwen as the vision model) gave four direct completions, one completion after guidance and one halt after
guidance. One of those completions is excluded because the spoken ``Thank you.'' was resolved to the previous
command's plan, which shows that the dialogue state carried over between runs. These runs differ in command,
planner and scene, so they are not a controlled comparison, and the final batch reached its targets with Gemma as
the vision model, so the 8~October failures were not a permanent limit of that model.

